\begin{abstract}
Protein-to-genome alignment---mapping protein sequences to their encoding genomic loci---is a foundational task in comparative genomics. While existing tools achieve near-perfect accuracy on one-to-one orthologs (exon F1 $>$ 0.87), we show they systematically fail on paralogs, the majority of gene families, with the best single tool reaching only 0.17--0.40 primary exon F1 across four vertebrate species. Through systematic causal diagnosis, we decompose this failure into two independent root causes: (1) a cross-species identity ceiling---human protein queries against non-human paralog loci cap at 0.64--0.71 F1 even given the correct locus (0.22--0.30 F1 loss); and (2) locus confusion---tools select the wrong paralog copy 51--89\% of the time (0.22--0.38 F1 loss). The diagnosis reveals a provable ceiling for annotation-free approaches: without species-matched sequence information, no alignment algorithm can resolve exon boundaries in diverged paralogs. Guided by this, we present TPRS---an alignment method that introduces annotation-derived auxiliary information to break through the ceiling. TPRS operates in two stages: Candidate Generation, which expands the search space along three dimensions---tool diversity (4 complementary A-stage tools), query identity (human + species-matched protein), and prior knowledge (Ensembl Compara co-ortholog recovery)---and Learned Joint Selection, where an XGBoost regressor scores every (locus, query-type) pair and selects the best by predicted exon F1. On paralogs, TPRS improves primary exon F1 by $+$0.36--0.46 over the best single alignment tool, recovering 87--90\% of the one-to-one ceiling in mammals and 71\% in zebrafish. Component ablation validates each design decision: species-matched protein is essential for primary F1 ($\Delta = -0.30$), learned routing is the largest single contributor ($\Delta = -0.39$), and each A-stage tool independently contributes to recall ($\Delta = -0.31$ per removal). TPRS demonstrates that for the paralog alignment problem---where sequence evidence alone is provably insufficient---annotation-derived priors transform an intractable task into one where high-quality spliced alignments can be reliably obtained.
\end{abstract}
